% GENERATED FILE. Edit canonical lessons, then run npm run generate:books.
% canonical-source-hash: fnv1a64:7c308068b9fc561a
% canonical-lessons: IT-C496-pannolino, IT-C496-passeggino, IT-C496-allattare, IT-C496-cullare, IT-C496-gattonare

\chapter{Il primo anno}
\label{ch:it-il-primo-anno}
\hlchaptermodality{496}

\begin{chapteropening}
\textbf{By the end of this chapter:} \emph{I can talk about looking after a baby.}

Complete the last lesson of chapter 496: I can talk about looking after a baby.
\end{chapteropening}

\section[il pannolino]{il pannolino — nappy, diaper}

\label{lesson:IT-C496-pannolino}



\begin{quote}
Before the new one: say the Italian for stepmother, then the Italian for half-brother.
\end{quote}

\subsection*{You'll want to know: il pannolino}
\textbf{il pannolino} — \textquotedblleft{}nappy, diaper\textquotedblright{}.

\begin{quote}
Devo cambiare il pannolino al bambino.
\end{quote}

\begin{quote}
\emph{I have to change the baby's nappy.}
\end{quote}

Built on \emph{panno} \textquotedblleft{}cloth\textquotedblright{}.

\begin{cousinweb}
\emph{pannolino} is a little form built on \emph{panno} \textquotedblleft{}cloth\textquotedblright{}, inherited from Latin \emph{pannus} \textquotedblleft{}a cloth, a rag\textquotedblright{} (Lewis \& Short).
\end{cousinweb}

\begin{grammarlens}[title={What you've built}]
One more word for telling, explaining and describing at B1.
\end{grammarlens}

\subsection*{Your turn}
\begin{itemize}
\raggedright
  \item \emph{Say it:} \emph{il pannolino}
  \item \emph{Say it:} \emph{il pannolino}, once more
  \item \emph{Say it:} a sentence of your own with \emph{il pannolino}
  \item \emph{Recall:} say the Italian for stepmother, then the Italian for half-brother, then say \emph{il pannolino} again
\end{itemize}

\begin{tcolorbox}[breakable,colback=teal!4,colframe=teal!35!black,title={Before you move on}]
What does il pannolino mean? (\textquotedblleft{}Nappy, diaper\textquotedblright{}.) Say it once more.
\end{tcolorbox}
\section[il passeggino]{il passeggino — pushchair, stroller}

\label{lesson:IT-C496-passeggino}



\begin{quote}
Before the new one: say the Italian for half-brother, then the Italian for nappy.
\end{quote}

\subsection*{You'll want to know: il passeggino}
\textbf{il passeggino} — \textquotedblleft{}pushchair, stroller\textquotedblright{}.

\begin{quote}
Porto la bambina al parco con il passeggino.
\end{quote}

\begin{quote}
\emph{I take the little girl to the park in the pushchair.}
\end{quote}

Built on \emph{passeggiare} \textquotedblleft{}to stroll\textquotedblright{}. A pram for a newborn is \emph{la carrozzina}.

\begin{cousinweb}
\emph{passeggino} is built on \emph{passeggiare} \textquotedblleft{}to go for a walk\textquotedblright{}, from \emph{passo} \textquotedblleft{}step\textquotedblright{}, inherited from Latin \emph{passus} \textquotedblleft{}a step\textquotedblright{} (Lewis \& Short).
\end{cousinweb}

\begin{grammarlens}[title={What you've built}]
One more word for telling, explaining and describing at B1.
\end{grammarlens}

\subsection*{Your turn}
\begin{itemize}
\raggedright
  \item \emph{Say it:} \emph{il passeggino}
  \item \emph{Say it:} \emph{il passeggino}, once more
  \item \emph{Say it:} a sentence of your own with \emph{il passeggino}
  \item \emph{Recall:} say the Italian for half-brother, then the Italian for nappy, then say \emph{il passeggino} again
\end{itemize}

\begin{tcolorbox}[breakable,colback=teal!4,colframe=teal!35!black,title={Before you move on}]
What does il passeggino mean? (\textquotedblleft{}Pushchair, stroller\textquotedblright{}.) Say it once more.
\end{tcolorbox}
\section[allattare]{allattare — to breastfeed, to nurse}

\label{lesson:IT-C496-allattare}



\begin{quote}
Before the new one: say the Italian for nappy, then the Italian for pushchair.
\end{quote}

\subsection*{You'll want to know: allattare}
\textbf{allattare} — \textquotedblleft{}to breastfeed, to nurse\textquotedblright{}.

\begin{quote}
Ha allattato il figlio per sei mesi.
\end{quote}

\begin{quote}
\emph{She breastfed her son for six months.}
\end{quote}

Built on \emph{latte} \textquotedblleft{}milk\textquotedblright{}.

\begin{cousinweb}
\emph{allattare} is usually traced to Late Latin \emph{allactāre} \textquotedblleft{}to suckle\textquotedblright{}, built on \emph{lac, lactis} \textquotedblleft{}milk\textquotedblright{} (Lewis \& Short).
\end{cousinweb}

\begin{grammarlens}[title={What you've built}]
One more word for telling, explaining and describing at B1.
\end{grammarlens}

\subsection*{Your turn}
\begin{itemize}
\raggedright
  \item \emph{Say it:} \emph{allattare}
  \item \emph{Say it:} \emph{allattare}, once more
  \item \emph{Say it:} a sentence of your own with \emph{allattare}
  \item \emph{Recall:} say the Italian for nappy, then the Italian for pushchair, then say \emph{allattare} again
\end{itemize}

\begin{tcolorbox}[breakable,colback=teal!4,colframe=teal!35!black,title={Before you move on}]
What does allattare mean? (\textquotedblleft{}To breastfeed, to nurse\textquotedblright{}.) Say it once more.
\end{tcolorbox}
\section[cullare]{cullare — to rock (a baby)}

\label{lesson:IT-C496-cullare}



\begin{quote}
Before the new one: say the Italian for pushchair, then the Italian for to breastfeed.
\end{quote}

\subsection*{You'll want to know: cullare}
\textbf{cullare} — \textquotedblleft{}to rock (a baby)\textquotedblright{}.

\begin{quote}
La mamma culla il bambino finché si addormenta.
\end{quote}

\begin{quote}
\emph{The mother rocks the baby until he falls asleep.}
\end{quote}

Built on \emph{la culla} \textquotedblleft{}cradle\textquotedblright{}.

\begin{cousinweb}
\emph{cullare} is built on \emph{culla} \textquotedblleft{}cradle\textquotedblright{}, inherited from Vulgar Latin \emph{cunla} (EtymDB), a shortened little form of \emph{cūnae} \textquotedblleft{}cradle\textquotedblright{} (Lewis \& Short has \emph{cūnulae} \textquotedblleft{}a little cradle\textquotedblright{}).
\end{cousinweb}

\begin{grammarlens}[title={What you've built}]
One more word for telling, explaining and describing at B1.
\end{grammarlens}

\subsection*{Your turn}
\begin{itemize}
\raggedright
  \item \emph{Say it:} \emph{cullare}
  \item \emph{Say it:} \emph{cullare}, once more
  \item \emph{Say it:} a sentence of your own with \emph{cullare}
  \item \emph{Recall:} say the Italian for pushchair, then the Italian for to breastfeed, then say \emph{cullare} again
\end{itemize}

\begin{tcolorbox}[breakable,colback=teal!4,colframe=teal!35!black,title={Before you move on}]
What does cullare mean? (\textquotedblleft{}To rock (a baby)\textquotedblright{}.) Say it once more.
\end{tcolorbox}
\section[gattonare]{gattonare — to crawl (on all fours)}

\label{lesson:IT-C496-gattonare}



\begin{quote}
Before the new one: say the Italian for to breastfeed, then the Italian for to rock.
\end{quote}

\subsection*{You'll want to know: gattonare}
\textbf{gattonare} — \textquotedblleft{}to crawl (on all fours)\textquotedblright{}.

\begin{quote}
A otto mesi mio figlio gattonava per tutta la casa.
\end{quote}

\begin{quote}
\emph{At eight months my son was crawling all over the house.}
\end{quote}

Built on \emph{gatto}: moving like a cat.

\begin{cousinweb}
\emph{gattonare} is built on \emph{gatto} \textquotedblleft{}cat\textquotedblright{}, inherited from Late Latin \emph{cattus} (EtymDB): to move on all fours like a cat.
\end{cousinweb}

\begin{grammarlens}[title={What you've built}]
One more word for telling, explaining and describing at B1.
\end{grammarlens}

\subsection*{Your turn}
\begin{itemize}
\raggedright
  \item \emph{Say it:} \emph{gattonare}
  \item \emph{Say it:} \emph{gattonare}, once more
  \item \emph{Say it:} a sentence of your own with \emph{gattonare}
  \item \emph{Recall:} say the Italian for to breastfeed, then the Italian for to rock, then say \emph{gattonare} again
\end{itemize}

\begin{tcolorbox}[breakable,colback=teal!4,colframe=teal!35!black,title={Before you move on}]
What does gattonare mean? (\textquotedblleft{}To crawl (on all fours)\textquotedblright{}.) Say it once more.
\end{tcolorbox}
